\clearpage
\SourceLeaf{194}
\thispagestyle{empty}
\pdfbookmark[0]{Allgemeine Auflösung der Aufgabe, die Theile einer gegebnen Fläche auf einer andern abzubilden}{abbildung}
\vspace*{.35in}
\begin{center}
{\large ALLGEMEINE AUFLÖSUNG DER AUFGABE}\par\bigskip
{\Large DIE THEILE EINER GEGEBNEN FLÄCHE}\par\bigskip
AUF EINER ANDERN GEGEBNEN FLÄCHE SO ABZUBILDEN\par\bigskip
DASS DIE ABBILDUNG DEM ABGEBILDETEN\par\bigskip
IN DEN KLEINSTEN THEILEN ÄHNLICH WIRD\par\vspace{28pt}
VON\par\vspace{28pt}
{\large CARL FRIEDRICH GAUSS}\par\vspace{30pt}
{\small ALS BEANTWORTUNG DER VON DER KÖNIGLICHEN SOCIETÄT DER WISSENSCHAFTEN\par
IN COPENHAGEN FÜR MDCCCXXII AUFGEGEBNEN PREISFRAGE.}\par\bigskip
\end{center}
\begin{flushright}‘Ab his via sternitur ad maiora.’\end{flushright}
\vfill
\begin{center}
\rule{.8\linewidth}{.3pt}\par\smallskip
Astronomische Abhandlungen herausgegeben von H. C. \textsc{Schumacher}.\par
Drittes Heft. Altona 1825.\par\smallskip
\rule{.8\linewidth}{.3pt}
\end{center}
\clearpage
